\begin{enumerate}
    \item[\textbf{Problema 2.1.}] Justifique a afirmação feita logo acima de que ``nove (9) gráficos diferentes, cada um contendo três (3) curvas'' são necessários para relacionar força e velocidade [3].

    \item[\textbf{Problema 2.2.}] Encontre e compare a densidade de massa e a viscosidade da manteiga de amendoim, mel e água [3].

    \item[\textbf{Problema 2.3.}] Qual é a constante na eq. (2.12)? Como você sabe disso? [4].

    \item[\textbf{Problema 2.4.}] Aplicar o método de base à eq. (2.2) para o período do pêndulo [4].

    \item[\textbf{Problema 2.5.}] Efectuar o método de base para a eq. (2.13) e mostrar que a massa de um corpo em queda não afecta a sua velocidade de descida [5].

    \item[\textbf{Problema 2.6.}] Execute a última etapa do método básico para eqs. (2.20) usando a primeira massa e mostrando-a produz uma forma equivalente à eq. (2.21) [5].

    \item[\textbf{Problema 2.7.}] Escreva o teorema de Buckingham Pi como uma série de etapas, análogas às etapas descritas no método básico na p. 23 [6].

    \item[\textbf{Problema 2.8.}] Confirmar a eq. (2.28) aplicando o método básico ao problema do misturador [6].

    \item[\textbf{Problema 2.9.}] Confirme a eq. (2.30) aplicando o resto do teorema de Buckingham Pi ao problema do pêndulo [6].

    \item[\textbf{Problema 2.10.}] Confirme a eq. (2.32) aplicando o resto do teorema de Buckingham Pi à formulação revisada do problema do pêndulo [6].

    \item[\textbf{Problema 2.11.}] Aplique o teorema de Buckingham Pi à revolução de dois corpos no espaço, começando com a equação funcional (2.16) [6, 7].

    \item[\textbf{Problema 2.12.}] O que acontece quando o teorema de Pi é aplicado ao problema dos dois corpos, mas começando agora com a equação funcional (2.14)? [7].

    \item[\textbf{Problema 2.13.}] Considere uma corda de comprimento $l$ que conecta uma rocha de massa $m$ a um ponto fixo enquanto a rocha gira em um círculo na velocidade $v$. Use o método básico de análise dimensional para mostrar que a tensão $T$ na corda é determinada pelo grupo adimensional [1, 8]:
    \begin{equation}
        \frac{TL}{m v^2} = \text{constante}
    \end{equation}

    \item[\textbf{Problema 2.14.}] Aplicar o teorema de Buckingham Pi para confirmar a análise do Problema 2.13 [8].

    \item[\textbf{Problema 2.15.}] A velocidade do som num gás, $c$, depende da pressão do gás $p$ e da densidade mássica do gás $\rho$. Use a análise dimensional para determinar como $c$, $p$ e $\rho$ estão relacionados [8].

    \item[\textbf{Problema 2.16.}] Um agrupamento adimensional chamado número de Weber, $We$, é usado em mecânica dos fluidos para relacionar a tensão superficial de um fluido em fluxo, $\sigma$, que tem dimensões de força/comprimento, com a velocidade do fluido, $v$, densidade, $\rho$, e um comprimento característico, $l$. Use a análise dimensional para encontrar esse número [8, 9].

    \item[\textbf{Problema 2.17.}] Um pêndulo oscila em um fluido viscoso. Quantos grupos são necessários para relacionar as variáveis usuais do pêndulo com a viscosidade do fluido, $\mu$, a densidade da massa do fluido, $\rho$ e o diâmetro $d$ do pêndulo? Encontre esses grupos [9].

    \item[\textbf{Problema 2.18.}] Pensa-se que o caudal volumétrico $Q$ do fluido através de um tubo depende da queda de pressão por unidade de comprimento, $\Delta p / l$, do diâmetro do tubo, $d$, e da viscosidade do fluido, $\mu$. Use o método básico de análise dimensional para determinar a relação [10]:
    \begin{equation}
        Q = (\text{constante}) \left(\frac{d^4}{\mu}\right) \left(\frac{\Delta p}{l}\right)
    \end{equation}

    \item[\textbf{Problema 2.19.}] Aplicar o teorema de Buckingham Pi para confirmar a análise do Problema 2.18 [10].

    \item[\textbf{Problema 2.20.}] Quando o fluxo em um tubo com uma parede interna áspera (talvez devido ao acúmulo de depósitos minerais) é considerado, várias variáveis devem ser consideradas, incluindo a velocidade do fluido $v$, sua densidade $\rho$ e viscosidade $\mu$, e o comprimento do tubo $l$ e diâmetro $d$. A variação média $e$ do raio do tubo pode ser tomada como uma medida da rugosidade da superfície interna do tubo. Determine os grupos adimensionais necessários para determinar como a queda de pressão $p$ depende dessas variáveis [10, 11].

    \item[\textbf{Problema 2.21.}] Use a análise dimensional para determinar como a velocidade do som no aço depende do módulo de elasticidade, $E$ e da densidade do aço, $\rho$. (O módulo de elasticidade do aço é, aproximadamente e em unidades britânicas, $30 \times 10^6$ psi.) [11, 12].

    \item[\textbf{Problema 2.22.}] A flexibilidade (a deflexão por unidade de carga) ou a complacência $C$ de uma viga com secção quadrada $d \times d$ depende do comprimento $l$ da viga, da sua altura e largura e do módulo de elasticidade $E$ do seu material. Use o método básico de análise dimensional para mostrar que [12]:
    \begin{equation}
        CEd = F_{CEd}\left(\frac{l}{d}\right)
    \end{equation}

    \item[\textbf{Problema 2.23.}] Experimentos foram conduzidos para determinar a forma específica da função $F_{CEd}(l/d)$ encontrada no Problema 2.22. Nesses experimentos, descobriu-se que um gráfico de $\log_{10} (CEd)$ contra $\log_{10} (l / d)$ tem uma inclinação de $3$ e uma interceptação na escala $\log_{10} (CEd)$ de $-0,60$. Mostre que a deflexão sob uma carga $P$ pode ser dada em termos do segundo momento da área $I$ da seção transversal ($I = d^4/12$) como [2, 12]:
    \begin{equation}
        \text{deflexão} = \text{carga} \times \text{conformidade} = P \times C = \frac{Pl^3}{48EI}
    \end{equation}

\end{enumerate}